% Independent mathematical derivation. Prefix CTP. Requires amsmath,amssymb,hyperref.
\section{The exact Cayley--Toeplitz receiver of the original entire zero trace}
\label{sec:cayley-toeplitz-full-receiver}

\paragraph{Objects and provenance.}
The entire function is the original Rodgers--Tao function
$H(Z)=H_0(Z)=\xi(1/2+iZ/2)/8$, with its original factor $1/8$,
coordinate $Z$, and positive theta kernel. The source is Brad Rodgers
and Terence Tao, \href{https://arxiv.org/abs/1801.05914v5}{\emph{The
de Bruijn--Newman constant is non-negative}}, the displayed definitions
of $\Phi$ and $H_t$ and the identity $H_0=\xi/8$.
The complete local-jet and Laurent receiver used here is (ALM1)--(ALM32),
\emph{A finite rational Laurent witness for every off-real heat zero};
the retained positive block and its exact residual are (RWR18)--(RWR27),
\emph{Independent reproduction of the supplied rational witness}.
The Toeplitz comparison source is Alain Connes and Walter D. van
Suijlekom, \href{https://arxiv.org/abs/2511.23257v1}{\emph{Quadratic
Forms, Real Zeros and Echoes of the Spectral Action}}, original
\path{Araki-final-oct25.tex},
the introduction's Corollary \texttt{corcar} and Section 2,
Proposition \texttt{toeplitzcase}, read directly in source lines
751--949. Their statement is the Carath\'eodory--Fej\'er theorem's
corank-one consequence, with its positivity and rank hypotheses.
The present calculation connects that receiver to Toeplitz forms;
it does not identify a finite Toeplitz eigenvector with a theta kernel.
All implications and maps used for this connection are proved below.
The support carrier is the original foundation retained under the
byline The Clankers,
\href{https://github.com/KokunoYumeto/zeta-function-research-reader/blob/a2851f9eb352e7e4376bc629de86fa4ed9c1c434/workbenches/splitzero-tandem/foundations/split-support-geometry-v11/split_support_geometry_arithmetic_curve_v11.tex}{\emph{Split Support Geometry}, v11, Definition \texttt{def:lattice-split}}.

\subsection{The actual function, every zero, and the complete rational domain}

Keep, without rescaling,
\begin{equation}\tag{CTP1}
 H(Z)=\int_0^\infty\Phi(u)\cos(Zu)\,du,\qquad
 \Phi(u)=\sum_{n\ge1}(2\pi^2n^4e^{9u}-3\pi n^2e^{5u})
                         e^{-\pi n^2e^{4u}}.
\end{equation}
Every summand is positive for $u\ge0$. The bound
$0<\Phi(u)\le C e^{9u-(\pi/2)e^{4u}}$, with
$C=2\pi^2\sum_{n\ge1}n^4e^{-(\pi/2)n^2}<\infty$, follows by
discarding the negative summand and using
$n^2e^{4u}\ge(n^2+e^{4u})/2$. Thus $H$ is entire, even, and real on
the real axis, and differentiation under its integral is valid on
every compact set. In particular
\begin{equation}\tag{CTP2}
 H(iy)=\int_0^\infty\Phi(u)\cosh(yu)\,du>0
 \quad(y\in\mathbb R).
\end{equation}
Write $\mathcal Z$ for its distinct zeros and $m_w$ for the full
multiplicity of $w\in\mathcal Z$. The estimates and paired
factorization in (ALM2)--(ALM4) give
\begin{equation}\tag{CTP3}
 \sum_{w\in\mathcal Z}\frac{m_w}{|w|^2}<\infty,\qquad
 H(Z)=H(0)\prod_{\{w,-w\}}(1-Z^2/w^2)^{m_w},\qquad
 \ell(Z):=\frac{H'(Z)}{H(Z)}
   =\sum_{\{w,-w\}}\frac{2m_wZ}{Z^2-w^2}.
\end{equation}
Here the product takes one representative of each sign pair, whereas
all sums below labeled $\mathcal Z$ take both roots and all
multiplicities. For reference, the bound on $\Phi$ gives
$\log\max_{|Z|\le R}|H(Z)|=O(R\log(2+R))$ by maximizing
$Ru-(\pi/8)e^{4u}$. Jensen's formula gives
$n(R)=O(R\log(2+R))$, and dyadic summation gives the first statement.
Hadamard factorization, paired under $w\mapsto-w$, then gives the
displayed product: the remaining exponential is of degree at most
one and its linear term is zero by evenness. The logarithmic
derivative follows locally uniformly off the zeros. These growth
and product facts matter: evenness and the zero set alone would not
exclude multiplying $H$ by $e^{Z^2}$ and changing $\ell$.
There are infinitely many distinct zeros, since a finite product
would be a polynomial, whereas every coefficient
$(-1)^j\int u^{2j}\Phi(u)du/(2j)!$ is nonzero.

Fix a real $b>0$. By (CTP2), $0,ib,-ib$ are not zeros. Let
\begin{equation}\tag{CTP4}
 \mathscr A_b=\{F\in\mathbb C(Z):
   \text{all finite poles of }F\text{ belong to }\{0,ib,-ib\},
   \ F(Z)=O(Z^{-1})\text{ at infinity}\}.
\end{equation}
This is a complex vector space, preserved by
$F^\dagger(Z)=\overline{F(\bar Z)}$. It contains the original Laurent
space $\mathscr R=\{Z^{-1}p(Z^{-2}):p\in\mathbb C[X]\}$.
Its full local-jet map and pairing are
\begin{equation}\tag{CTP5}
\begin{split}
 j_b(F)&=([F]_w)_{w\in\mathcal Z}
   \in\prod_{w\in\mathcal Z}\mathcal O_{\mathbb C,w}/(H),\\
 \mathcal T(F,G)&=\sum_{w\in\mathcal Z}m_w
                  \overline{F(\bar w)}G(w)
  =\sum_{w\in\mathcal Z}\operatorname{Res}_{Z=w}
          F^\dagger(Z)G(Z)\ell(Z)\,dZ.
\end{split}
\end{equation}
The series is absolutely convergent by (CTP3), since its summands
are $O(m_w|w|^{-2})$ outside a finite set. Every local algebra retains
its original length $m_w$: if $H=(Z-w)^mU$ with $U(w)\ne0$, then
$\ell=m/(Z-w)+U'/U$, giving the displayed residue. Only the resulting
trace kills nilpotents. The jet map itself retains them. A rational
function vanishing in all these jets vanishes at infinitely many
distinct points and is zero, so $j_b$ is injective. The infinite
pairing is defined on $\mathscr A_b$ or its specified injective image,
not on an unrestricted product of local algebras. Conjugation of the
sum and the substitution $w\mapsto\bar w$ prove that $\mathcal T$
is Hermitian.

The original arithmetic residue factors remain exactly
$s=1/2+iZ/2$, $g(s)=16H(-2i(s-1/2))$,
$ds/g=(i/32)dZ/H$, and $\alpha^*g'=-32iH'$. Their product is one.
Consequently (CTP5) is the same Jacobian residue contraction as
(ALM7)--(ALM8), extended to the indicated rational domain.

\subsection{A finite partial-fraction rule for the entire zero trace}

Suppose $r\in\mathbb C(Z)$ has poles only in $\{0,ib,-ib\}$ and is
$O(Z^{-2})$. Its unique partial fraction expansion is
\begin{equation}\tag{CTP6}
 r(Z)=\sum_{a\in\{0,ib,-ib\}}\sum_{n=1}^{N_a}
                  A_{a,n}(Z-a)^{-n},\qquad
 \sum_a A_{a,1}=0.
\end{equation}
The last equality is its vanishing coefficient of $Z^{-1}$ at
infinity. Every coefficient is an exact finite derivative:
$A_{a,n}=\left[(\partial_Z^{N_a-n}/(N_a-n)!)
((Z-a)^{N_a}r(Z))\right]_{Z=a}$ when $N_a$ is the pole order.
Then
\begin{equation}\tag{CTP7}
 \sum_{w\in\mathcal Z}m_w r(w)
  =-\sum_{a\in\{0,ib,-ib\}}\sum_{n=1}^{N_a}
          \frac{A_{a,n}}{(n-1)!}\ell^{(n-1)}(a).
\end{equation}
For $n\ge2$, differentiating the paired, locally uniform series in
(CTP3) gives
$\sum_w m_w(w-a)^{-n}=-\ell^{(n-1)}(a)/(n-1)!$.
For the simple-pole terms, use $\sum_a A_{a,1}=0$ before summation:
$\sum_a A_{a,1}(w-a)^{-1}
=\sum_a A_{a,1}((w-a)^{-1}-w^{-1})$.
Each difference is $O(|w|^{-2})$ and its sum is $-\ell(a)$ by the
paired product and $\ell(0)=0$. This proves (CTP7) without assigning
an absolutely convergent meaning to an isolated $\sum_w1/w$.
Apply (CTP7) to $r=F^\dagger G$ to compute every finite matrix of
(CTP5) from finitely many derivatives of the actual $H$ at
$0,ib,-ib$. It is an exact finite evaluation rule with all zeros
present, not a truncation of the zero multiset.

\subsection{Cayley coordinates, the Toeplitz matrix, and its exact finite map}

Set
\begin{equation}\tag{CTP8}
 C_b(Z)=\frac{Z-ib}{Z+ib},\qquad
 F_j(Z)=\frac{C_b(Z)^j}{Z+ib}\quad(j\in\mathbb Z),\qquad
 c_k(b)=\sum_{w\in\mathcal Z}
               \frac{m_w}{w^2+b^2}C_b(w)^k\quad(k\in\mathbb Z).
\end{equation}
The values $C_b(w)$ are nonzero and finite. They tend to one as
$|w|\to\infty$. Thus, for every fixed integer $k$, $C_b(w)^k$ is
bounded over the root set and $|w^2+b^2|^{-1}=O(|w|^{-2})$.
This proves absolute convergence of every $c_k$ and of every
pairing of the $F_j$.

The involution has the precise, necessary orientation
\begin{equation}\tag{CTP9}
 C_b^\dagger=C_b^{-1},\qquad
 F_i^\dagger(Z)=\frac{C_b(Z)^{-i}}{Z-ib}=F_{-i-1}(Z),\qquad
 F_i^\dagger F_j=\frac{C_b^{j-i}}{Z^2+b^2}.
\end{equation}
For $d\ge1$ define the complex-linear injection
\begin{equation}\tag{CTP10}
 W_{b,d}:\mathbb C^d\longrightarrow\mathscr A_b,\quad
 a\longmapsto\sum_{j=0}^{d-1}a_jF_j,\qquad
 T_d(b)=(c_{j-i}(b))_{0\le i,j<d}.
\end{equation}
Equations (CTP5) and (CTP9) prove, with no scalar factor missing,
\begin{equation}\tag{CTP11}
 \mathcal T(W_{b,d}a,W_{b,d}v)=a^*T_d(b)v.
\end{equation}
The map is injective because $C_b$ is a nonconstant rational
coordinate. The identities $C_b(-w)=C_b(w)^{-1}$ and
$C_b(-\bar w)=\overline{C_b(w)}$, together with the unchanged
multiplicities and weights, give
\begin{equation}\tag{CTP12}
 c_{-k}=c_k=\overline{c_k}.
\end{equation}
Thus the exact matrix is real symmetric and Toeplitz. This conclusion
uses both evenness and reality of the original $H$.

The finite coefficient map also has an explicit rational inverse.
If $p(U)=\sum_{j<d}a_jU^j$, write
\begin{equation}\tag{CTP13}
 W_{b,d}a(Z)=\frac{q(Z)}{(Z+ib)^d},\qquad
 q(Z)=\sum_{j=0}^{d-1}a_j(Z-ib)^j(Z+ib)^{d-j-1}.
\end{equation}
This maps all polynomials of degree less than $d$ bijectively to
all such $q$; the inverse is
\begin{equation}\tag{CTP14}
 p(U)=\frac{(1-U)^{d-1}}{(2ib)^{d-1}}
                  q\!\left(ib\frac{1+U}{1-U}\right).
\end{equation}
The right side is a polynomial of degree at most $d-1$ because
$\deg q<d$. Substituting $U=C_b(Z)$ and
$1-U=2ib/(Z+ib)$ verifies both compositions exactly. For rational
$b>0$ these maps have coefficients in $\mathbb Q(i)$.
The union of the images is precisely the rational functions
vanishing at infinity whose only possible pole is $-ib$.

\subsection{An exact resolvent from the original logarithmic derivative}

For $z\ne1$ put $v_b(z)=-ib(1+z)/(1-z)$, and let
$R_b=\sup_{w\in\mathcal Z}|C_b(w)|<\infty$. The geometric series,
dominated by the absolutely summable weights for $|z|R_b<1$, gives
\begin{equation}\tag{CTP15}
\begin{split}
 G_b(z):=\sum_{k\ge0}c_kz^k
 &=\sum_{w\in\mathcal Z}
   \frac{m_w}{(w^2+b^2)(1-zC_b(w))}\\
 &=\frac{\ell(v_b(z))-\ell(ib)}{2ib}.
\end{split}
\end{equation}
Indeed
\[
 \frac1{(w^2+b^2)(1-zC_b(w))}
 =\frac1{2ib}\left(\frac1{w-ib}-\frac1{w-v_b(z)}\right).
\]
The difference is summable and (CTP3) sums it to the stated
logarithmic derivatives. The expression is meromorphic on
$\mathbb C\setminus\{1\}$. Alternatively its root series converges
locally uniformly off $1$ and the indicated poles: outside a finite
root set, $C_b(w)$ is uniformly close to one and the denominators
stay uniformly away from zero on such a compact set. Its only
finite poles away from $1$ are
\begin{equation}\tag{CTP16}
 z_w=C_b(w)^{-1}=\frac{w+ib}{w-ib},\qquad
 \operatorname{Res}_{z=z_w}G_b(z)=-\frac{m_w}{(w-ib)^2}\ne0.
\end{equation}
The Cayley map is injective, so distinct roots give distinct poles.
For a chosen $z_w\ne1$, choose a closed disk about $z_w$ disjoint
from $1$ and all other poles. Such a disk exists: an infinite
sequence of distinct roots has modulus tending to infinity, and
its Cayley pole images tend to $1$; the remaining finitely many
images are distinct. On this disk the series with the one $w$
term removed is locally uniformly convergent and holomorphic,
by the preceding denominator bound. It cannot cancel the pole.
The multiplicity enters the nonzero residue and cannot disappear.
For a direct sign check, $v_b'(z)=-2ib/(1-z)^2$, so the residue in
the last expression of (CTP15) is
$m_w(1-z_w)^2/(4b^2)=-m_w/(w-ib)^2$.

Putting $z=0$ and using $\ell(-ib)=-\ell(ib)$ gives
\begin{equation}\tag{CTP17}
 c_0=-\frac{H'(ib)}{ibH(ib)}
 =\frac{\int_0^\infty u\Phi(u)\sinh(bu)du}
           {b\int_0^\infty\Phi(u)\cosh(bu)du}>0.
\end{equation}
The strict inequality uses the actual positive kernel in (CTP1),
not merely the fact that $H(ib)>0$.
For every $k\ge1$, expand $v_b(z)=-ib-2ibz/(1-z)$ near zero.
The coefficient of $z^k$ in $(z/(1-z))^r$ is
$\binom{k-1}{r-1}$ for $1\le r\le k$. Since
$\ell^{(r)}(-ib)=(-1)^{r+1}\ell^{(r)}(ib)$, one obtains the finite
formula
\begin{equation}\tag{CTP18}
 c_k=-\sum_{r=1}^k
       \binom{k-1}{r-1}\frac{(2ib)^{r-1}}{r!}\ell^{(r)}(ib).
\end{equation}
All derivatives are of the original $H$; its original mass has not
been reassigned. For example, $c_1=-\ell'(ib)$ equals
\begin{equation}\tag{CTP19}
 c_1=\frac{\int u^2\Phi(u)\cosh(bu)du}{\int\Phi(u)\cosh(bu)du}
 -\left(\frac{\int u\Phi(u)\sinh(bu)du}
                  {\int\Phi(u)\cosh(bu)du}\right)^2>0.
\end{equation}
Cauchy--Schwarz bounds the subtracted numerator by
$(\int\Phi\cosh)(\int u^2\Phi\sinh^2/\cosh)$, which is strictly
less than the first product because
$\cosh-\sinh^2/\cosh=1/\cosh>0$. This proves the last inequality.
It does not prove $c_1\le c_0$ or any later Toeplitz inequality.

\subsection{Completeness: every nonreal zero yields a two-coefficient test}

If all the roots are real, (CTP8)--(CTP11) give the positive finite
measure and its exact Toeplitz form
\begin{equation}\tag{CTP20}
 \mu_b=\sum_{w\in\mathcal Z}\frac{m_w}{w^2+b^2}\delta_{C_b(w)},
 \quad \mu_b(\mathbb T)=c_0,\quad
 a^*T_da=\int_{\mathbb T}|p(U)|^2d\mu_b(U)>0\quad(a\ne0).
\end{equation}
Here $\mathbb T=\{U:|U|=1\}$. Strictness follows from infinitely
many distinct atoms and the fact that a nonzero polynomial has only
finitely many zeros. No division by $c_0$ has been made.

If a nonreal root exists, conjugation supplies one with
$\operatorname{Im}w<0$. Directly,
\begin{equation}\tag{CTP21}
 |C_b(w)|^2=
 \frac{(\operatorname{Re}w)^2+(\operatorname{Im}w-b)^2}
      {(\operatorname{Re}w)^2+(\operatorname{Im}w+b)^2}>1.
\end{equation}
Since $C_b(w)\to1$ along unbounded roots, $R_b>1$ is then a
maximum, attained in a finite set. Equations (CTP15)--(CTP16) show
that the Taylor series at zero has radius exactly $1/R_b<1$:
there is no pole at smaller modulus, and a pole at that modulus
has a nonzero residue. Therefore
\begin{equation}\tag{CTP22}
 \limsup_{k\to\infty}|c_k|^{1/k}=R_b>1.
\end{equation}
In particular there is a finite $k\ge1$ with $|c_k|>c_0$.
Because $c_k$ is real, set $\varepsilon_k=\operatorname{sgn}c_k$
and keep the two coefficients exactly equal to $1,-\varepsilon_k$:
\begin{equation}\tag{CTP23}
 f_{b,k}(Z)=\frac{1-\varepsilon_kC_b(Z)^k}{Z+ib},\qquad
 \mathcal T(f_{b,k},f_{b,k})=2(c_0-|c_k|)<0.
\end{equation}
This is a finite integer coefficient vector in the original
$T_{k+1}$, with zero entries in all positions except $0,k$.
For rational $b$ the function itself has coefficients in
$\mathbb Q(i)$. The series giving its trace still includes every
root and every multiplicity.

For any fixed $b>0$ we have consequently proved the equivalence
\begin{equation}\tag{CTP24}
\begin{aligned}
 &\text{all zeros of the actual }H\text{ are real}\\
 &\quad\Longleftrightarrow T_d(b)\succ0\text{ for every }d\ge1\\
 &\quad\Longleftrightarrow T_d(b)\succeq0\text{ for every }d\ge1\\
 &\quad\Longleftrightarrow |c_k(b)|\le c_0(b)\text{ for every }k\ge1.
\end{aligned}
\end{equation}
For the final implication use the contrapositive (CTP22); the
forward implication also follows from (CTP20). Positive
semidefiniteness implies the last bound by the principal block on
indices $0,k$ in $T_{k+1}$. These are completeness statements for
the constructed receiver. No assertion that their infinite
positivity condition holds has been made. In the real-root case
the Taylor radius is exactly one: the poles $z_w$ lie on the unit
circle, are individually nonzero, and accumulate only at one.
Thus in both cases the growth identity
$\limsup|c_k|^{1/k}=R_b$ holds, with $R_b=1$ in the real-root case.

\subsection{The exact relation to all Laurent moments, at every value of the parameter}

Retain the reciprocal nodes and their entire multiplicity fibers
from (ALM10): $\Lambda=\{w^{-2}:w\in\mathcal Z\}$,
$M_\lambda=2m_w$ when $\lambda=w^{-2}$,
$V=\sum_\lambda M_\lambda|\lambda|<\infty$, and
$R=\max_\lambda|\lambda|$. Then, for $k\ge0$,
\begin{equation}\tag{CTP25}
\begin{split}
 P_k(Y)&=\sum_{l=0}^k(-1)^l\binom{2k}{2l}Y^l,\\
 c_k&=\sum_{\lambda\in\Lambda}M_\lambda\lambda
                 \frac{P_k(b^2\lambda)}{(1+b^2\lambda)^{k+1}}.
\end{split}
\end{equation}
To verify the identity, average the $w$ and $-w$ summands of
(CTP8), using unchanged multiplicities. Their even rational factor is
\[
 \frac{C_b(w)^k+C_b(w)^{-k}}{2(w^2+b^2)}
 =\frac{\sum_{l=0}^k(-1)^l\binom{2k}{2l}b^{2l}w^{2k-2l}}
             {(w^2+b^2)^{k+1}}.
\]
This is precisely the function of $\lambda$ in (CTP25).
Its denominator never vanishes on the actual node set because
$H(ib)\ne0$. The aggregation loses neither member of a sign pair.

When $b^2R<1$, the exact convergent expansion in the original
Laurent moments $S_n=\sum_wm_ww^{-2n}$ is
\begin{equation}\tag{CTP26}
 c_k=\sum_{n\ge0}(-1)^nb^{2n}S_{n+1}
       \sum_{l=0}^{\min(k,n)}\binom{2k}{2l}\binom{k+n-l}{n-l}.
\end{equation}
Expand $(1+Y)^{-k-1}$ by the binomial series and multiply by
$P_k(Y)$. The displayed inner sum is its coefficient with sign
$(-1)^n$. Absolute convergence and interchange with the root sum
follow from $|S_{n+1}|\le VR^n$ and a polynomial bound in $n$ for
the inner sum, with the fixed ratio $b^2R<1$.
This restricted series is not used outside its stated disk.

For every $b>0$, even when that series does not converge, the
rational function in (CTP25) is a controlled limit of actual
Laurent polynomials. Here is a constructive proof that also fixes
the omitted infinite tail. Let
$r_k(X)=P_k(b^2X)/(1+b^2X)^{k+1}$. Choose
$0<r_0<r_1<1/b^2$ with no node of modulus $r_0$ and let
$S=\{\lambda\in\Lambda:|\lambda|>r_0\}$, a finite set.
The Taylor series of $r_k$ converges uniformly on $|X|\le r_0$;
for any $\epsilon>0$ take a Taylor polynomial $Q$ whose error there
is less than $\epsilon/2$. For $\lambda\in S$ define
$L_\lambda(X)=\prod_{\mu\in S\setminus\{\lambda\}}
(X-\mu)/(\lambda-\mu)$ and
$B_\lambda=\prod_{\mu\ne\lambda}(r_0+|\mu|)/|\lambda-\mu|$.
Choose a finite integer $N$ such that
\begin{equation}\tag{CTP27}
 \sum_{\lambda\in S}|r_k(\lambda)-Q(\lambda)|B_\lambda
                   (r_0/|\lambda|)^N<\epsilon/2,
 \quad
 Q_\epsilon(X)=Q(X)+\sum_{\lambda\in S}
 [r_k(\lambda)-Q(\lambda)](X/\lambda)^NL_\lambda(X).
\end{equation}
If $S$ is empty the correction sum is zero. Each ratio in a
nonempty sum is strictly less than one, so the choice exists.
The correction has the prescribed value at each node in $S$,
and on the whole remaining disk its absolute size is less than
$\epsilon/2$. Consequently
\begin{equation}\tag{CTP28}
 \sup_{\lambda\in\Lambda}|Q_\epsilon(\lambda)-r_k(\lambda)|
       <\epsilon,\qquad
 \left|\sum_\lambda M_\lambda\lambda Q_\epsilon(\lambda)-c_k\right|
       \le V\epsilon.
\end{equation}
Conjugation pairs make $Q_\epsilon$ real when $Q$ is taken real.
Every sum on the left is a finite linear combination of the
original $S_n$. Thus (CTP25)--(CTP28) give an exact, globally
valid bridge from the Laurent receiver, with a proved full-tail
bound. They do not delete a cutoff-dependent boundary term.

There is no nonzero finite identification of the two original
test spaces by equality of their full jets: a Laurent test has
its only possible pole at zero, whereas a Cayley test in the
union of (CTP10) has its only possible pole at $-ib$. Equality
in every jet implies rational equality by the injectivity in
(CTP5); such a common function would have no finite pole and
vanish at infinity, hence be zero. Their precise relation is
the common injective rational receiver (CTP4)--(CTP7), the
moment maps (CTP25)--(CTP28), and the following exact residual.

\subsection{The original positive block and the Cayley residual}

Keep the certified block $A=K_{64}(0)\succ0$ of (RWR18), in its
original coordinates, and the original injection
$E:\mathbb C^{64}\to\mathscr R$, $Ea=\sum_{h=0}^{63}a_hZ^{-2h-1}$.
The positivity and its interval certificate are already part of
(RWR1)--(RWR17); no new numerical assertion is used here.
For $F\in\mathscr A_b$ set $r(F)_h=\mathcal T(Ee_h,F)$ and define
\begin{equation}\tag{CTP29}
 \Pi_bF=F-EA^{-1}r(F),\qquad
 \mathcal T(Ea,\Pi_bF)=0,\qquad \Pi_b^2=\Pi_b.
\end{equation}
These equalities follow from $r(Ea)=Aa$. The kernel of $\Pi_b$
is exactly $E\mathbb C^{64}$, and its image is the orthogonal
complement of that space in $\mathscr A_b$. For the original
Laurent subspace this restricts exactly to (RWR25).
For $0\le j<d$ set
\begin{equation}\tag{CTP30}
 B_{hj}=\mathcal T(Z^{-2h-1},F_j),\qquad
 D_d(b)=T_d(b)-B^*A^{-1}B,\qquad
 \widetilde W_{b,d}=\Pi_bW_{b,d}.
\end{equation}
The entries of $B$ have the exact finite derivative evaluation
(CTP7) with
$r(Z)=Z^{-2h-1}(Z-ib)^j/(Z+ib)^{j+1}$.
All constants, signs, and poles are thus specified, including
the full root sum. Direct expansion proves
\begin{equation}\tag{CTP31}
\begin{split}
 \mathcal T(\widetilde W a,\widetilde W v)&=a^*D_d(b)v,\\
 \mathcal T(E\alpha+W a,E\alpha+W a)
  &=(\alpha+A^{-1}Ba)^*A(\alpha+A^{-1}Ba)+a^*D_d(b)a.
\end{split}
\end{equation}
The map $\widetilde W$ is injective: a vector in its kernel
would give a common nonzero Laurent and Cayley function,
which the preceding pole argument excludes. The block coordinate
change is exactly $(\alpha,a)\mapsto(\alpha+A^{-1}Ba,a)$,
with inverse $(\beta,a)\mapsto(\beta-A^{-1}Ba,a)$.
Every negative test (CTP23) has
$a^*D_da=a^*T_da-(Ba)^*A^{-1}(Ba)<0$.
Conversely, positivity of every $D_d$ implies positivity of every
$T_d$ by (CTP30), hence reality of all roots by (CTP24).
If all roots are real, (CTP5) is positive on every nonzero rational
test, so every $D_d$ is strictly positive by injectivity.
This proves a complete alternative residual criterion, in the
same original trace, with no assertion about its unresolved sign.
On the direct sum of the original Laurent and Cayley test spaces,
(CTP29) also retains every mixed matrix entry; replacing that direct
sum by two independently positive blocks would discard those entries.

\subsection{The entire support lattice and all connecting maps}

Retain the original arbitrary nontrivial bounded distributive
lattice $L$. For every vector space $V$ used above keep
\begin{equation}\tag{CTP32}
\begin{gathered}
 G_L(V)=\{(0,\lambda):\lambda\in L\}\cup(V\times\{1_L\}),\\
 (v,\lambda)+(w,\mu)=(v+w,\lambda\vee\mu),\quad
 (a,\eta)(v,\lambda)=(av,\eta\wedge\lambda),\\
 \tau=(0,0_L),\qquad e=(0,1_L).
\end{gathered}
\end{equation}
For every complex-linear map $A:V\to W$ in (CTP5), (CTP10),
(CTP13)--(CTP14), and (CTP29)--(CTP31), the exact lift
$A^\uparrow(v,\lambda)=(Av,\lambda)$ is well defined,
additive, and $G_L(\mathbb C)$-linear. Nonzero output amplitude
requires nonzero input amplitude and hence top support; the
operation checks are direct substitution and lattice
distributivity. The pairing lift is
\begin{equation}\tag{CTP33}
 \mathcal T^L((F,\lambda),(G,\mu))
       =(\mathcal T(F,G),\lambda\wedge\mu),\qquad
 \mathcal T^L(W^\uparrow(a,\lambda),W^\uparrow(v,\mu))
       =(a^*T_dv,\lambda\wedge\mu).
\end{equation}
It is well defined because a nonzero pairing requires both
amplitudes nonzero. Sesquilinearity, including the unchanged
involution $(a,\eta)^*=(\bar a,\eta)$, follows from the scalar
pairing and distributivity of meet over join. The same formula
applies to $D_d$, $B$, and the original Laurent residual.

Retain the programme's idempotent semifield $K$ with zero,
$K_L=K\otimes_{\mathbb B}L$, and $\iota_L(\lambda)=1_K\otimes\lambda$.
The full comparison and its commuting maps are
\begin{equation}\tag{CTP34}
 \Phi_V(v,\lambda)=(v,\iota_L(\lambda)),\qquad
 (A\times I)\Phi_V=\Phi_WA^\uparrow.
\end{equation}
The nonzero character $\chi:K\to\mathbb B$ is a semiring map:
in an idempotent semifield a sum is zero only when both inputs
are zero, and a product is zero only when one input is zero.
Therefore $\rho_L=\chi\otimes I_L$ satisfies
$\rho_L\iota_L=I_L$. This proves that the comparison retains
every original support label. Equations (CTP33)--(CTP34)
commute with the pairing, with meet of the two input labels.
Every canceled amplitude at top support remains $e$ and is
not changed to $\tau$. The finite coordinate decompositions
use $G_L(\mathbb C^{64}\oplus\mathbb C^d)$ with its one retained
label, not a product with independently chosen labels.

In particular the negative test (CTP23), when present, has
the unchanged image
\begin{equation}\tag{CTP35}
 (2(c_0-|c_k|),1_L)\xmapsto{\Phi_{\mathbb R}}
       (2(c_0-|c_k|),1_K\otimes1_L).
\end{equation}
The original tau geometry, every zero multiplicity, the complete
arithmetic residue factor, and all finite or infinite moment
indices have survived the construction. The mathematics above
constructs and proves the receiver and its precise detection
criterion; it does not evaluate the sign of all its forms.

\subsection{The exact Li and Weil maps, including the original factors}

The coefficient criterion originates in Xian-Jin Li,
\href{https://doi.org/10.1006/jnth.1997.2137}{\emph{The positivity of
a sequence of numbers and the Riemann hypothesis}} (1997).
Enrico Bombieri and Jeffrey C. Lagarias developed its zero-multiset
and Weil-functional relations in
\href{https://doi.org/10.1006/jnth.1999.2392}{\emph{Complements to
Li's criterion for the Riemann hypothesis}} (1999).
The original author LaTeX actually consulted here is Jeffrey C.
Lagarias, \href{https://arxiv.org/abs/math/0404394v4}{\emph{Li
Coefficients for Automorphic $L$-Functions}},
\path{math-0404394v4.tex}, lines 197--359 and 1170--1424:
source equations \texttt{103}, \texttt{101}, \texttt{102a},
\texttt{301}, \texttt{303}, \texttt{304}, and Theorem
\texttt{th31}. The author explicitly credits J. B. Keiper's 1992
precursor coefficients in lines 263--274. The following map
identifies the present objects with those coefficients and test
functions; no historical novelty for the coefficient criterion is
asserted.

For arbitrary $b>0$ define, in a neighborhood of $z=0$, the unique
holomorphic logarithm with value zero at zero by
\begin{equation}\tag{CTP36}
 \log\frac{H(v_b(z))}{H(ib)}
     =\sum_{n\ge1}\frac{\Lambda_n(b)}n z^n,
 \qquad \Lambda_0(b)=0,
 \qquad \Lambda_{-n}(b)=\Lambda_n(b).
\end{equation}
This is legitimate because $H(v_b(0))=H(-ib)=H(ib)>0$.
There is no change to $H$ or its mass: the displayed ratio merely
specifies the logarithm's constant. For each sign pair of roots,
the exact factor identity is
\[
 \frac{1-v_b(z)^2/w^2}{1+b^2/w^2}
 =\frac{(1-zC_b(w))(1-zC_b(w)^{-1})}{(1-z)^2}.
\]
Logarithmically differentiate the locally uniformly convergent
paired product, or take coefficients in its logarithm on a small
zero-free disk. This gives
\begin{equation}\tag{CTP37}
 \Lambda_n(b)=\sum_{\{w,-w\}}m_w
                 (2-C_b(w)^n-C_b(w)^{-n})\qquad(n\ge1).
\end{equation}
The summand is $O(m_w|w|^{-2})$: for fixed $n$ the first-order
terms in $C_b(w)-1=O(1/w)$ cancel between the two powers.
Thus the paired sum is absolutely convergent with every
multiplicity retained.
Differentiating (CTP36) and using
$v_b'=-2ib/(1-z)^2$, (CTP15), and
$\ell(ib)=-ibc_0$, gives
\begin{equation}\tag{CTP38}\begin{gathered}
 (1-z)^2\sum_{n\ge1}\Lambda_n(b)z^{n-1}
           =4b^2G_b(z)-2b^2c_0,\\
 \qquad
 \Lambda_1(b)=2b^2c_0,\quad
 \Lambda_{k+1}(b)-2\Lambda_k(b)+\Lambda_{k-1}(b)=4b^2c_k.\end{gathered}
\end{equation}
The last equality is valid for $k\ge1$; with the displayed even
extension and $\Lambda_0=0$ it also holds for $k=0$.
Its inverse, obtained by summing these second differences, is
\begin{equation}\tag{CTP39}
 \Lambda_n(b)=2b^2nc_0+4b^2\sum_{k=1}^{n-1}(n-k)c_k
 \quad(n\ge1).
\end{equation}
Indeed both sides have the same second differences and the same
values at $n=0,1$. These are finite, mutually inverse linear maps
on each required coefficient segment, with no lost constant term.

At $b=1$, the original arithmetic coordinate gives
\begin{equation}\tag{CTP40}
 s(v_1(z))=\frac1{1-z},\qquad
 \frac{H(v_1(z))}{H(i)}
     =\frac{\xi(1/(1-z))}{\xi(0)},\qquad
 \xi(0)=\xi(1).
\end{equation}
Lagarias's equations \texttt{102a} and \texttt{103c} therefore
identify $\Lambda_n(1)$ with the classical Li coefficient
$\lambda_n$ (his index $-n$ has the same value).
The exact specialization is
\begin{equation}\tag{CTP41}
 c_0(1)=\frac{\lambda_1}{2},\qquad
 c_k(1)=\frac{\lambda_{k+1}-2\lambda_k+\lambda_{k-1}}4
 \quad(k\ge1).
\end{equation}
In particular the complete two-coefficient detector (CTP24)
is exactly the bound
$|\lambda_{k+1}-2\lambda_k+\lambda_{k-1}|\le2\lambda_1$
for every $k\ge1$, in these unchanged conventions. No proof of
that infinite family of bounds is claimed.

The test-space identification is stronger than an agreement of
coefficient sequences. Let
$G_n(s)=1-(1-1/s)^n$ be Lagarias's Li test function, for
$n\in\mathbb Z$, and let $\mathscr L$ be his Li class: rational
functions vanishing at infinity with poles contained in $\{0,1\}$.
Under $\alpha(Z)=1/2+iZ/2$ one has exactly
\begin{equation}\tag{CTP42}
 C_1(Z)=\frac{s}{s-1},\qquad
 F_j(Z)=\alpha^*\left(\frac{G_{-j-1}-G_{-j}}{2i}\right)(Z)
 \quad(j\in\mathbb Z).
\end{equation}
The second identity follows from
$G_{-j-1}-G_{-j}=C_1^j(1-C_1)=C_1^j/(1-s)$ and
$Z+i=2i(1-s)$. For $m>0$ the inverse finite sums are
\begin{equation}\tag{CTP43}
 \alpha^*G_{-m}=2i\sum_{j=0}^{m-1}F_j,
 \qquad \alpha^*G_m=-2i\sum_{j=-m}^{-1}F_j.
\end{equation}
They telescope directly. Equations (CTP13)--(CTP14) for the two
pole locations also show that the span of all integer-indexed
$F_j$ is exactly the pullback of $\mathscr L$.
The additional Laurent pole $Z=0$ corresponds to $s=1/2$; it
is retained in $\mathscr A_1$ and is not claimed to belong to
Lagarias's test class. The full local-jet trace is its proved
common receiving space.

Lagarias's source pairing, equation \texttt{301}, is linear in
its first entry. The programme pairing is conjugate-linear
in its first entry. With this difference explicitly retained,
\begin{equation}\tag{CTP44}
 \mathcal T(\alpha^*f,\alpha^*g)
       =\langle g,f\rangle_{\mathcal W},\qquad
 \langle G_n,G_m\rangle_{\mathcal W}
       =\lambda_n+\lambda_{-m}-\lambda_{n-m}.
\end{equation}
For the first equality use $\alpha(\bar w)=1-\overline{\alpha(w)}$
in every summand and preserve $m_w$. For the second, the finite
identity
$G_n(s)G_{-m}(s)=G_n(s)+G_{-m}(s)-G_{n-m}(s)$ and
$\overline{G_m(1-\bar s)}=G_{-m}(s)$ prove the formula with
the convergent symmetric zero sums used by the source.
Their combined product is absolutely summable. This also gives
a direct verification of the displayed source theorem without
depending on any typographical intermediate expansion in its
proof. Substituting (CTP42) into (CTP44) gives
\begin{equation}\tag{CTP45}
 (T_d(1))_{ij}=\frac14
   (\lambda_{j-i+1}-2\lambda_{j-i}+\lambda_{j-i-1}),
 \qquad \lambda_{-n}=\lambda_n.
\end{equation}
The factor is $1/4$, the product of $1/(2i)$ with its conjugate.
Thus this Toeplitz form is exactly the finite-difference pullback
of the classical Li--Weil Gram form, with complete multiplicities.
All coordinate, finite-sum, and coefficient maps
(CTP38)--(CTP45) are fixed complex-linear maps on their stated
spaces, so their explicit full-$L$ lifts are (CTP32)--(CTP34).
For the map reversing the two pairing entries, the labels are
reversed as well; their meet is unchanged by commutativity.

\subsection{The precise finite extremal defect in the source comparison}

For any particular $d$, let $\mu_d$ be the actual least eigenvalue
of the actual Hermitian matrix $T_d(b)$. The spectral theorem
gives the exact positive semidefinite Toeplitz matrix
\begin{equation}\tag{CTP46}
 \widehat T_d=T_d(b)-\mu_d I_d\succeq0,
 \qquad \ker\widehat T_d=\ker(T_d(b)-\mu_d I_d).
\end{equation}
This construction does not assert $\mu_d\ge0$ or simplicity.
Its changed coefficient is exactly $c_0-\mu_d$; every
nonconstant coefficient $c_k$, $|k|<d$, remains $c_k$.
The Carath\'eodory--Fej\'er corank-one consequence in Connes--van
Suijlekom, Corollary \texttt{corcar} and Proposition
\texttt{toeplitzcase}, concerns this positive matrix when its
kernel has dimension one. Its polynomial is the coefficient
polynomial of a vector in that kernel. Equations (CTP13)--(CTP14)
give the exact Cayley rational function of that polynomial;
they provide no identity of this finite polynomial with $H$.

The full defect is an actual functional, not an unspecified
change of matrix. Define on Laurent polynomials in $U$
\begin{equation}\tag{CTP47}
 \mathscr L_b(P)=\sum_{w\in\mathcal Z}
            \frac{m_w}{w^2+b^2}P(C_b(w)),\qquad
 h(P)=\frac1{2\pi}\int_0^{2\pi}P(e^{i\theta})d\theta,
 \qquad \widehat{\mathscr L}_{b,d}=\mathscr L_b-\mu_d h.
\end{equation}
Every sum converges because a fixed Laurent polynomial is
bounded on the Cayley image of the roots. With
$P^*(U)=\overline{P(1/\bar U)}$, (CTP9) gives
$\mathscr L_b(P^*Q)=\mathcal T(P(C_b)/(Z+ib),Q(C_b)/(Z+ib))$.
Furthermore $h(U^k)=\delta_{k0}$ and
$h(P^*P)=\sum_j|a_j|^2$ for $P=\sum a_jU^j$.
For the finite restriction take exactly
$P(U)=\sum_{j=0}^{d-1}a_jU^j$ and
$Q(U)=\sum_{j=0}^{d-1}v_jU^j$, with $a,v\in\mathbb C^d$.
Then
\begin{equation}\tag{CTP48}
 \widehat{\mathscr L}_{b,d}(P^*Q)=a^*\widehat T_dv,
 \qquad
 \mathcal T(Wa,Wa)=\mu_d\sum_j|a_j|^2
 \quad(a\in\ker\widehat T_d).
\end{equation}
Every factor in this equality is retained, including the
$1/(2\pi)$ in the unit-mass Haar functional. It is not a
rescaling of the original zero-trace mass $c_0$.
The analytic generating function of the changed moments is
exactly $G_b(z)-\mu_d$. Subtracting a constant leaves every
pole (CTP16) and its residue unchanged. Thus finite extremal
positivity does not remove the all-zero defect that (CTP22)
detects. Equations (CTP46)--(CTP48) are the exact connecting
map and its defect required to apply the source's Toeplitz
geometry to this receiver.

Finally $\mathscr L_b$, $h$, and
$\widehat{\mathscr L}_{b,d}$ are complex-linear maps on the same
Laurent-polynomial space. Their full lifts send
$(P,\lambda)$ respectively to
$(\mathscr L_b(P),\lambda)$, $(h(P),\lambda)$ and
$(\mathscr L_b(P)-\mu_dh(P),\lambda)$.
Cancellation in the difference retains $\lambda$ exactly;
at top support it is $e$, never $\tau$. These maps commute
with $\Phi$ by direct substitution in (CTP34), and their
Hermitian forms carry the meet of the two input labels.
